\documentclass[../main.tex]{subfiles}

\begin{document}
\begin{CJK*}{UTF8}{gbsn}

\section*{Exercise 3}

Consider split conformal prediction with $\mathcal{Y} = \mathbb{R}$,
two given models $\hat{f}_i: \mathcal{X} \to \mathbb{R}$
for $i = 1, 2$, and the score function 

\begin{equation*}
  s(x,y) = \min_{i = 1, 2} \abs{y - \hat{f}_i(x)}.
\end{equation*}
Derive an explicit expression for $\mathcal{C}(x)$
in terms of $\hat{f}_1(x)$, $\hat{f}_2(x)$, and the quantile $\hat{q}$ of the scores on the validation set.
What is the shape of the prediction when $\hat{q}$ is small or when 
it is large?
Briefly explain why this might be attractive if the conditional 
distribution of $Y$ given $X = x$ is bimodal.

\paragraph{Solution}
By definition, it is immediately clear that

\begin{equation*}
    \mathcal{C}(x) = \{y \in \mathcal{Y} \mid s(x,y) \leqslant \hat{q}\}
    = [\hat{f}_1(x) - \hat{q}, \hat{f}_1(x) + \hat{q}] \cup [\hat{f}_2(x) - \hat{q}, \hat{f}_2(x) + \hat{q}].
\end{equation*}
If $\hat{f}_1(x)=\hat{f}_2(x)$, then $\mathcal{C}(x)$ is an interval 
of length $2\hat{q}$ centered at $\hat{f}_1(x) = \hat{f}_2(x)$.
If we assume that the two models are different at $x$,
that is, 
$\hat{f}_1(x) \neq \hat{f}_2(x)$, then $\mathcal{C}(x)$ is the union of two intervals
when $\hat{q}$ is small enough, and it is a single interval when $\hat{q}$ is large enough.
If the conditional 
distribution of $Y$ given $X = x$ is bimodal,
then the score might be of interest 
because $\hat{f}_1(x)$ and $\hat{f}_2(x)$
might capture the two modes of the distribution 
and the predicted interval will be centred at both modes. 
////
\end{CJK*}
\end{document}